\chapter{退出路径：IPO、并购与二级转让}

\section{退出是投资 thesis 的一部分}

进入时就应讨论\textbf{最可能的退出通道}：战略并购、财务买家、IPO、 secondary sale（卖老股）或 recap。不同通道对持股比例、治理清洁度、审计年限与合规记录要求不同；临退出才补洞往往代价高昂。

\section{并购退出}

战略买家通常愿付\textbf{协同溢价}，但谈判周期与反垄断审查不确定。财务买家（另一家 PE）依赖杠杆与再次出售，尽调更偏财务模型。拍卖流程（broad auction vs targeted）影响价格与泄密风险；顾问费用与管理层激励（management equity rollover）需在条款中明确。

\section{IPO}

公开市场提供流动性与品牌，但面临\textbf{锁定期、持续披露成本与市场波动}。上市地选择（本地 vs 离岸）涉及估值倍数、监管环境与外汇安排。GP 需与承销商、公司律师协调\textbf{注册权}与减持节奏。

\section{二级份额转让}

基金到期或 LP 流动性需求可能触发\textbf{基金二级交易}（LP 份额转让）或\textbf{投资组合二级出售}（continuation fund）。定价依赖 NAV 与尽职调查；利益冲突（GP 同时代表买卖双方）须严格程序化并披露。

\section{基金期限与延期}

LPA 规定的基金到期后，未退出资产可通过\textbf{延期}（通常需一定比例 LP 同意）、\textbf{捆绑出售}或\textbf{接续基金}处理。延期费率与 GP 激励常重新谈判；透明沟通可减少纠纷。
